\documentclass[10pt,a4paper]{article}
\usepackage{../../styles/neet_style}

\begin{document}

% Title Header
\begin{center}
  {\Huge \textbf{NEET UG 2027: Master Biology Chapter Notes}}\\[6pt]
  {\Large \textbf{Reproduction in Plants \& Humans, Reproductive Health \& Biotechnology}}\\[4pt]
  {\small \textbf{Prepared, Compiled \& Authored by Satya Prakash} $\cdot$ Strictly Calibrated to Rationalized NCERT \& NMC Scope}
  \vspace{4pt}
  \hrule height 1.5pt
\end{center}

\pagestyle{fancy}
\fancyhf{}
\fancyhead[L]{\small \textbf{Reproduction \& Biotechnology Master Notes}}
\fancyhead[R]{\small Satya Prakash}
\fancyfoot[L]{\small \textbf{NEET UG 2027 Prep Series}}
\fancyfoot[C]{\thepage}
\fancyfoot[R]{\small \textbf{Biology Chapter Notes}}

\section{Sexual Reproduction in Flowering Plants (Angiosperms)}

\begin{conceptbox}[Microsporogenesis \& Male Gametophyte (Pollen Grain)]
\begin{itemize}
  \item \textbf{Anther Structure:} Typically bilobed, dithecous (two thecae per lobe), and tetrasporangiate (four microsporangia).
  \item \textbf{Wall Layers (from outside to inside):}
  \begin{enumerate}
    \item \textbf{Epidermis:} Protective single outer layer.
    \item \textbf{Endothecium:} Layer with fibrous hygroscopic $\alpha$-cellulose bands; assists in anther dehiscence.
    \item \textbf{Middle Layers:} $1-3$ transient ephemeral layers; degenerate to provide nourishment.
    \item \textbf{Tapetum:} Innermost nutritive layer; cells possess dense cytoplasm and are generally \textbf{multinucleate / polyploid}; secretes \textbf{sporopollenin}, \textbf{callase}, and pollenkitt.
  \end{enumerate}
  \item \textbf{Pollen Grain Wall (Sporoderm):}
  \begin{itemize}
    \item \textbf{Exine:} Hard outer layer composed of \textbf{Sporopollenin} (the most resistant organic substance known; withstands high temperatures, strong acids, and alkali; no enzyme that degrades sporopollenin is known). Has apertures called \textbf{Germ Pores} where sporopollenin is absent.
    \item \textbf{Intine:} Thin, continuous inner wall composed of \textbf{cellulose and pectin}.
  \end{itemize}
  \item \textbf{Pollen Shedding Stage:} In $>60\%$ of angiosperms, pollen is shed at the \textbf{2-celled stage} (large vegetative cell with irregular nucleus $+$ small spindle-shaped generative cell floating in vegetative cytoplasm). In $<40\%$, generative cell divides mitotically to form 2 male gametes before shedding (\textbf{3-celled stage}).
\end{itemize}
\end{conceptbox}

\begin{conceptbox}[Megasporogenesis \& The Monosporic Polygonum-Type Embryo Sac]
\begin{itemize}
  \item \textbf{Ovule (Integumented Megasporangium):} Attached to placenta by \textbf{funicle}; point of attachment to ovule body is \textbf{hilum}. Massive central tissue is \textbf{nucellus}.
  \item \textbf{Megasporogenesis:} Megaspore Mother Cell (MMC) undergoes meiosis to produce a linear tetrad of 4 haploid megaspores. 3 micropylar megaspores degenerate; only \textbf{1 chalazal megaspore remains functional} (\textbf{Monosporic development}).
  \item \textbf{Embryo Sac Organization (7-celled, 8-nucleate at maturity):}
  \begin{itemize}
    \item \textbf{Micropylar End (Egg Apparatus - 3 cells):} 1 central \textbf{Egg cell} $+$ 2 flanking \textbf{Synergids} possessing \textbf{Filiform apparatus} (finger-like wall projections that guide pollen tube entry).
    \item \textbf{Chalazal End (Antipodals - 3 cells):} Vegetative nutritive cells; degenerate after fertilization.
    \item \textbf{Central Cell (1 large cell):} Contains \textbf{2 Polar Nuclei} that fuse before fertilization to form a diploid \textbf{Secondary Nucleus} ($2n$).
  \end{itemize}
\end{itemize}
\end{conceptbox}

\begin{formulabox}[Double Fertilization, Endosperm \& Apomixis]
\textbf{1. Double Fertilization (Discovered by S.G. Nawaschin, 1898):}
Unique to angiosperms; involves two distinct nuclear fusion events:
\begin{align*}
  \textbf{Syngamy: } & \text{Male Gamete } (n) + \text{Egg } (n) \longrightarrow \mathbf{\text{Zygote } (2n)} \longrightarrow \text{Embryo}\\
  \textbf{Triple Fusion: } & \text{Male Gamete } (n) + \text{Secondary Nucleus } (2n) \longrightarrow \mathbf{\text{PEN } (3n)} \longrightarrow \text{Endosperm}\\
  \textbf{Total Nuclei: } & 1 \text{ (sperm)} + 1 \text{ (egg)} + 1 \text{ (sperm)} + 2 \text{ (polar nuclei)} = \mathbf{5\,\text{Nuclei}}
\end{align*}
\textbf{2. Endosperm Development:}
Triploid Primary Endosperm Nucleus (PEN) divides to form endosperm before the zygote divides (ensures nutrition for embryo).
\begin{itemize}
  \item \textbf{Free-Nuclear Endosperm:} PEN undergoes repeated free nuclear divisions without cell wall formation (e.g., coconut water is free-nuclear endosperm consisting of thousands of free nuclei; surrounding white kernel is cellular endosperm).
\end{itemize}
\textbf{3. Special Reproductive Modes:}
\begin{itemize}
  \item \textbf{Apomixis:} Form of asexual reproduction that \textbf{mimics sexual reproduction}, producing seeds without fertilization (e.g., \textit{Asteraceae} and grasses). Nucellar cells divide and protrude into embryo sac to form diploid embryos.
  \item \textbf{Polyembryony:} Occurrence of more than one embryo in a seed (common in \textit{Citrus} and mango).
\end{itemize}
\end{formulabox}

\section{Human Reproduction \& Reproductive Health}

\begin{conceptbox}[Male \& Female Reproductive Architecture]
\textbf{1. Male System:}
\begin{itemize}
  \item \textbf{Scrotum:} Keeps testes at $2-2.5^\circ\mathrm{C}$ lower than internal core body temperature, necessary for spermatogenesis.
  \item \textbf{Seminiferous Tubules:} Lined by (i) \textbf{Spermatogonia} (male germ cells) and (ii) \textbf{Sertoli cells} (nurse cells providing nutrition, secrete Androgen Binding Protein and Inhibin).
  \item \textbf{Interstitial / Leydig Cells:} Located in interstitial spaces outside tubules; synthesize and secrete androgens (\textbf{Testosterone}) under LH stimulation.
  \item \textbf{Accessory Glands:} Paired Seminal vesicles (contribute $60-70\%$ of semen, rich in \textbf{fructose, calcium}, and prostaglandins), single Prostate gland (milky alkaline fluid), paired Bulbourethral / Cowper's glands (secrete alkaline mucus for lubricating penis).
\end{itemize}
\textbf{2. Female System:}
\begin{itemize}
  \item \textbf{Ovaries:} Primary female sex organs producing ovum and steroid hormones (Estrogen, Progesterone).
  \item \textbf{Fallopian Tube (Oviduct, 10--12 cm):} Consists of (i) \textbf{Infundibulum} with finger-like \textbf{fimbriae} to collect ovum after ovulation, (ii) \textbf{Ampulla} (the wider part; \textbf{site of fertilization}), (iii) \textbf{Isthmus} (narrow lumen joining uterus).
  \item \textbf{Uterine Wall:} Outer thin \textbf{perimetrium}, middle thick smooth-muscle \textbf{myometrium} (exhibits strong contractions during parturition), and inner glandular \textbf{endometrium} (undergoes cyclic changes during menstrual cycle).
\end{itemize}
\end{conceptbox}

\begin{formulabox}[Gametogenesis: Spermatogenesis vs Oogenesis Timing]
\textbf{1. Spermatogenesis (Continuous from Puberty through Old Age):}
\begin{align*}
  \text{Spermatogonia } (2n) & \xrightarrow{\text{mitosis}} \text{Primary Spermatocyte } (2n) \xrightarrow{\text{Meiosis I}} 2 \times \text{Secondary Spermatocytes } (n)\\
  & \xrightarrow{\text{Meiosis II}} 4 \times \text{Spermatids } (n) \xrightarrow{\mathbf{Spermiogenesis}} 4 \times \mathbf{\text{Spermatozoa } (n)}
\end{align*}
\textbf{2. Oogenesis (Initiated during Fetal Life, Suspended until Puberty):}
\begin{align*}
  \text{Oogonia } (2n) & \xrightarrow{\text{mitosis}} \text{Primary Oocyte } (2n) \xrightarrow{\text{Arrested at Prophase I}} \text{At Birth}\\
  \text{At Puberty: } & \text{Completes Meiosis I} \longrightarrow \mathbf{\text{Secondary Oocyte } (n)} + \mathbf{1^{\mathrm{st}} \text{ Polar Body } (n)}\\
  \text{Fertilization: } & \text{Completes Meiosis II} \longrightarrow \mathbf{\text{Ovum } (n)} + \mathbf{2^{\mathrm{nd}} \text{ Polar Body } (n)}
\end{align*}
\end{formulabox}

\begin{conceptbox}[Menstrual Cycle Hormonal Regulation \& Implantation]
Standard 28-day cycle divided into four phases:
\begin{enumerate}
  \item \textbf{Menstrual Phase (Days 1--5):} Breakdown of endometrial lining and blood vessels due to sudden \textbf{withdrawal of progesterone}; average blood loss $= 50-100\,\mathrm{mL}$.
  \item \textbf{Follicular / Proliferative Phase (Days 6--13):} Pituitary \textbf{FSH} stimulates maturation of primary follicle into mature \textbf{Graafian follicle}; growing follicle secretes \textbf{Estrogen}, which regenerates endometrium.
  \item \textbf{Ovulatory Phase (Day 14):} Peak levels of both LH and FSH; rapid surge of LH (\textbf{LH Surge}) induces rupture of Graafian follicle and release of secondary oocyte (\textbf{Ovulation}).
  \item \textbf{Luteal / Secretory Phase (Days 15--28):} Ruptured follicle transforms into endocrine \textbf{Corpus Luteum}, which secretes massive amounts of \textbf{Progesterone} (essential for maintenance of endometrium). If fertilization fails, corpus luteum degenerates into \textbf{Corpus Albicans}, causing progesterone drop and menses.
\end{enumerate}
\textbf{Fertilization \& Pregnancy Endocrinology:}
\begin{itemize}
  \item \textbf{Site of Fertilization:} Ampullary-isthmic junction / \textbf{Ampulla} of Fallopian tube.
  \item \textbf{Cortical Reaction:} Sperm induces depolarization and hardening of \textbf{zona pellucida}, blocking polyspermy.
  \item \textbf{Implantation:} Trophoblast layer of \textbf{Blastocyst} attaches to uterine endometrium $\approx 7$ days after fertilization.
  \item \textbf{Placental Hormones exclusively in pregnancy:} Human Chorionic Gonadotropin (\textbf{hCG}), Human Placental Lactogen (\textbf{hPL}), and \textbf{Relaxin} (secreted by ovary in late pregnancy).
\end{itemize}
\end{conceptbox}

\begin{conceptbox}[Reproductive Health: Contraception, MTP \& ART Protocols]
\textbf{1. Contraceptive Methodologies:}
\begin{itemize}
  \item \textbf{Natural / Traditional:} Periodic abstinence (avoid coitus days $10-17$ of cycle), coitus interruptus, \textbf{lactational amenorrhea} (intense lactation prevents ovulation for max 6 months postpartum).
  \item \textbf{Barrier Methods:} Condoms (protect against STIs/AIDS), Diaphragms, cervical caps, vaults.
  \item \textbf{Intrauterine Devices (IUDs):}
  \begin{itemize}
    \item \textbf{Non-medicated IUDs:} \textit{Lippes loop} (increases phagocytosis of sperms).
    \item \textbf{Copper-releasing IUDs:} \textbf{CuT, Cu7, Multiload 375} ($\mathrm{Cu}^{2+}$ ions suppress sperm motility and fertilizing capacity).
    \item \textbf{Hormone-releasing IUDs:} \textbf{Progestasert, LNG-20} (make uterus unsuitable for implantation and cervix hostile to sperms).
  \end{itemize}
  \item \textbf{Oral Contraceptive Pills:} Daily combined progestogen-estrogen pills; or \textbf{Saheli} (developed by CDRI Lucknow; non-steroidal, once-a-week pill with high contraceptive value and negligible side effects).
  \item \textbf{Surgical / Permanent Methods:} \textbf{Vasectomy} (cutting/tying vas deferens in males) and \textbf{Tubectomy} (fallopian tubes in females); irreversible sterilization.
\end{itemize}
\textbf{2. Assisted Reproductive Technologies (ART):}
\begin{itemize}
  \item \textbf{IVF-ET (In Vitro Fertilization - Embryo Transfer):} Test-tube baby program.
  \begin{itemize}
    \item \textbf{ZIFT (Zygote Intra-Fallopian Transfer):} Zygote or early embryo \textbf{up to 8 blastomeres} transferred into Fallopian tube.
    \item \textbf{IUT (Intra-Uterine Transfer):} Embryos with \textbf{more than 8 blastomeres} transferred directly into Uterus.
  \end{itemize}
  \item \textbf{GIFT (Gamete Intra-Fallopian Transfer):} Transfer of ovum collected from donor into fallopian tube of female who cannot produce one.
  \item \textbf{ICSI (Intra-Cytoplasmic Sperm Injection):} Sperm directly injected into ovum cytoplasm in laboratory.
  \item \textbf{IUI (Intra-Uterine Insemination):} Semen collected from husband/donor introduced artificially into uterus.
\end{itemize}
\end{conceptbox}

\section{Biotechnology: Principles, Tools \& Modern Applications}

\begin{conceptbox}[Tools of Recombinant DNA Technology]
\textbf{1. Restriction Endonucleases ("Molecular Scissors"):}
Discovered by Arber, Smith, and Nathans (1970). Cut double-stranded DNA at specific palindromic recognition sequences.
\begin{itemize}
  \item \textbf{First Discovered:} \textit{Hind II} (cuts at specific 6 bp sequence).
  \item \textbf{Nomenclature ($Eco\mathrm{RI}$):} $E$ for genus \textit{Escherichia}, $co$ for species \textit{coli}, $\mathrm{R}$ for strain RY13, $\mathrm{I}$ for first enzyme isolated.
  \item \textbf{Recognition Sequence of $Eco\mathrm{RI}$:}
  \begin{equation}
    5'-\mathrm{G \downarrow A \ A \ T \ T \ C}-3' \qquad 3'-\mathrm{C \ T \ T \ A \ A \uparrow G}-5'
  \end{equation}
  Produces staggered single-stranded overhangs called \textbf{sticky ends}.
\end{itemize}
\textbf{2. Cloning Vectors (Features of Ideal Vector pBR322):}
\begin{itemize}
  \item \textbf{Origin of Replication ($Ori$):} Specific sequence where replication initiates; controls copy number of linked foreign DNA.
  \item \textbf{Selectable Markers:} Identify and select transformants and eliminate non-transformants. Example: genes conferring resistance to antibiotics (\textbf{$amp^R$ and $tet^R$}).
  \item \textbf{Insertional Inactivation:} Ligation of foreign DNA at \textit{BamHI} site of $tet^R$ gene inactivates tetracycline resistance $\implies$ recombinant grows on ampicillin, dies on tetracycline.
  \item \textbf{Blue-White Screening:} Foreign DNA inserted into gene coding for enzyme \textbf{$\beta$-galactosidase} ($\alpha$-complementation). Non-recombinants produce blue colonies with X-gal; recombinants produce \textbf{white colonies}.
\end{itemize}
\textbf{3. Competent Host Transformation Techniques:}
\begin{itemize}
  \item \textbf{Chemical Method:} Treating bacterial cells with divalent cation ($\mathrm{Ca}^{2+}$) followed by \textbf{Heat Shock} ($42^\circ\mathrm{C}$ for 2 minutes and ice plunge).
  \item \textbf{Microinjection:} Recombinant DNA directly injected into nucleus of animal cell.
  \item \textbf{Biolistics / Gene Gun:} Plant cells bombarded with high-velocity microparticles of \textbf{gold or tungsten} coated with DNA.
  \item \textbf{Disarmed Pathogen Vectors:} \textit{Agrobacterium tumefaciens} (Ti plasmid for dicots) and retroviruses for animal cells.
\end{itemize}
\end{conceptbox}

\begin{formulabox}[Polymerase Chain Reaction (PCR) \& Downstream Processing]
Invented by Kary Mullis in 1983. In vitro enzymatic amplification of DNA sequence by a billion-fold ($2^n$ copies):
\begin{align*}
  \textbf{Step 1: Denaturation: } & 94-96^\circ\mathrm{C} \quad (\text{Thermal separation of dsDNA strands by breaking H-bonds})\\
  \textbf{Step 2: Annealing: } & 50-60^\circ\mathrm{C} \quad (\text{Hybridization of two oligonucleotide primers to flanking } 3'\text{-ends})\\
  \textbf{Step 3: Extension: } & 72^\circ\mathrm{C} \quad (\text{Extension by } \mathbf{\text{Taq Polymerase}} \text{ from } \textit{Thermus aquaticus})
\end{align*}
\textbf{Bioreactors \& Downstream Processing:}
\begin{itemize}
  \item \textbf{Bioreactors ($100-1000\,\mathrm{L}$):} Provide optimal growth conditions (pH, temperature, oxygen, substrates, salts). \textbf{Stirred-tank} provides continuous mixing; \textbf{sparged stirred-tank} bubbles sterile air through broth.
  \item \textbf{Downstream Processing:} Separation and purification of biosynthesized product, followed by clinical trials, formulation with preservatives, and stringent quality control.
\end{itemize}
\end{formulabox}

\begin{ncertbox}[High-Yield Transgenic Applications \& Medical Biotechnology]
\begin{enumerate}[leftmargin=*]
  \item \textbf{Bt Cotton \& Pest Resistance:}
  Soil bacterium \textit{Bacillus thuringiensis} produces crystalline proteinaceous protoxins (\textbf{Cry proteins}). The inactive protoxin is converted into active toxin in the \textbf{alkaline pH of the insect midgut}, binding to epithelial cells, creating pores, causing cell lysis and insect death:
  \begin{equation}
    \mathbf{CryIAc \text{ and } CryIIAb} \text{ control Cotton Bollworms} \qquad \mathbf{CryIAb} \text{ controls Corn Borer}
  \end{equation}
  \item \textbf{RNA Interference (RNAi):}
  Cellular defense mechanism in all eukaryotes. Nematode \textit{Meloidogyne incognita} infects roots of tobacco plants. Using \textit{Agrobacterium} vectors, nematode-specific genes are introduced producing both sense and anti-sense RNA $\implies$ forms \textbf{double-stranded RNA (dsRNA)} $\implies$ silences specific mRNA of nematode.
  \item \textbf{Genetically Engineered Insulin (Humulin):}
  Human insulin consists of chain A (21 amino acids) and chain B (30 amino acids) linked by \textbf{disulphide bonds}. Synthesized as \textbf{prohormone containing an extra C-peptide} (33 amino acids), which is removed during maturation. In 1983, American company \textbf{Eli Lilly} produced two separate DNA sequences for chain A and chain B, introduced into \textit{E. coli}, and combined them by creating disulphide bonds.
  \item \textbf{Gene Therapy for ADA Deficiency:}
  First clinical gene therapy performed in 1990 on a 4-year-old girl with \textbf{Adenosine Deaminase (ADA) deficiency} (causes Severe Combined Immunodeficiency - SCID). Lymphocytes from patient's blood transduced with functional ADA cDNA using a retroviral vector and returned to patient. (Requires periodic infusion unless embryonic stem cells are engineered).
  \item \textbf{Transgenic Animals \& Biopiracy:}
  \begin{itemize}
    \item \textbf{First Transgenic Cow Rosie (1997):} Produced human protein-enriched milk ($2.4\,\mathrm{g/litre}$) containing \textbf{human alpha-lactalbumin}, nutritionally more balanced for human babies than natural cow milk.
    \item \textbf{Biopiracy Case:} In 1997, US company RiceTec was granted patent on Basmati rice cross-bred with semi-dwarf varieties, misappropriating traditional Indian farmers' strains.
  \end{itemize}
\end{enumerate}
\end{ncertbox}

\end{document}
